% ===== PRÉAMBULE CANONIQUE — à coller VERBATIM en tête de CHAQUE .tex =====
\documentclass[11pt,a4paper]{article}
\usepackage[utf8]{inputenc}\usepackage[T1]{fontenc}\usepackage{lmodern}
\usepackage[french]{babel}
\usepackage{amsmath,amssymb,mathtools}\usepackage{icomma}
\usepackage[version=4]{mhchem}          % \ce{...} : équations & formules
\usepackage{siunitx}\sisetup{locale=FR} % \qty{}{}, \si{}, \ang{}
\usepackage{chemfig}                    % \chemfig{...} : structures développées
\usepackage{xcolor}\usepackage{colortbl}
\usepackage{tikz}\usepackage{pgfplots}\pgfplotsset{compat=1.18}
\usepackage[most]{tcolorbox}\usepackage{enumitem}\usepackage{array,booktabs}
\usepackage[a4paper,margin=2.2cm]{geometry}
\usetikzlibrary{arrows.meta,calc,angles,quotes,positioning,patterns,backgrounds,shapes.geometric}
\definecolor{primary}{RGB}{222,49,99}\definecolor{primarydark}{RGB}{150,22,55}
\definecolor{defcol}{RGB}{0,114,178}\definecolor{methcol}{RGB}{52,92,148}
\definecolor{excol}{RGB}{0,135,90}\definecolor{attncol}{RGB}{200,40,55}
\tcbset{boxbase/.style={enhanced,breakable,boxrule=0.9pt,arc=2.5pt,
  left=3mm,right=3mm,top=2.3mm,bottom=2.3mm,fonttitle=\bfseries,coltitle=white}}
% --- boîtes TUTORIEL (noms EXACTS, ne pas renommer) ---
\newtcolorbox{cle}[1][]{boxbase,colback=defcol!6,colframe=defcol,
  title={\textsc{À retenir}\ifx\relax#1\relax\relax\else\textmd{ : \itshape #1}\fi}}
\newtcolorbox{methode}[1][]{boxbase,colback=methcol!7,colframe=methcol,sharp corners,
  title={\textsc{Méthode}\ifx\relax#1\relax\relax\else\textmd{ --- \itshape #1}\fi}}
\newtcolorbox{exemplebox}[1][]{boxbase,colback=excol!5,colframe=excol,
  title={\textsc{Exemple}\ifx\relax#1\relax\relax\else\textmd{ : \itshape #1}\fi}}
\newtcolorbox{piege}[1][]{boxbase,colback=attncol!6,colframe=attncol,
  title={\textsc{Piège}\ifx\relax#1\relax\relax\else\textmd{ : \itshape #1}\fi}}
\newtcolorbox{remarque}[1][]{boxbase,colback=black!4,colframe=black!45,
  title={\textsc{Remarque}\ifx\relax#1\relax\relax\else\textmd{ : \itshape #1}\fi}}
% --- boîtes EXERCICE / CORRIGÉ (noms EXACTS, auto-comptées) ---
\newtcolorbox[auto counter]{exo}[1][]{boxbase,colback=primary!4,colframe=primary,
  title={\textsc{Exercice}~\thetcbcounter\ifx\relax#1\relax\relax\else\textmd{ --- \itshape #1}\fi}}
\newtcolorbox[auto counter]{corrige}[1][]{boxbase,colback=excol!4,colframe=excol,
  title={\textsc{Corrigé}~\thetcbcounter\ifx\relax#1\relax\relax\else\textmd{ --- \itshape #1}\fi}}
\newcommand{\R}{\mathbb{R}}\newcommand{\N}{\mathbb{N}}\newcommand{\Z}{\mathbb{Z}}\newcommand{\ds}{\displaystyle}
% (un \usepackage{fancyhdr}+en-tête est possible ; il est ignoré côté web)
% =========================================================================
\begin{document}
\sisetup{per-mode=symbol}
\begin{corrige}[identifier le réactif limitant]
Le plus petit rapport $n/\nu$ désigne le limitant.
\begin{enumerate}[label=\alph*)]
  \item \ce{C} : $\frac{3}{1}=3$ ; \ce{O2} : $\frac{2}{1}=2$. Plus petit : \ce{O2} $\Rightarrow$
        \textbf{\ce{O2} limitant}.
  \item \ce{H2} : $\frac{5}{2}=2,5$ ; \ce{O2} : $\frac{2}{1}=2$. Plus petit : \ce{O2} $\Rightarrow$
        \textbf{\ce{O2} limitant}.
  \item \ce{N2} : $\frac{2}{1}=2$ ; \ce{H2} : $\frac{9}{3}=3$. Plus petit : \ce{N2} $\Rightarrow$
        \textbf{\ce{N2} limitant}.
\end{enumerate}
\end{corrige}

\begin{corrige}[le réactif en excès qui reste]
\begin{enumerate}[label=\alph*)]
  \item \ce{H2} : $\frac{5}{2}=2,5$ ; \ce{O2} : $\frac{2}{1}=2$. \textbf{\ce{O2} est limitant}.
  \item $n(\ce{H2})_{\text{consommée}}=n(\ce{O2})\times\dfrac{2}{1}=2\times2=\qty{4}{\mole}.$
  \item $n(\ce{H2})_{\text{restant}}=5-4=\qty{1}{\mole}.$
\end{enumerate}
\end{corrige}

\begin{corrige}[calculer le produit à partir du limitant]
\begin{enumerate}[label=\alph*)]
  \item \ce{Al} : $\frac{4}{2}=2$ ; \ce{Cl2} : $\frac{3}{3}=1$. \textbf{\ce{Cl2} est limitant}.
  \item $n(\ce{AlCl3})=n(\ce{Cl2})\times\dfrac{2}{3}=3\times\dfrac{2}{3}=\qty{2}{\mole}.$
  \item \ce{Al} consommé $=3\times\frac{2}{3}=\qty{2}{\mole}$, donc restant $=4-2=\qty{2}{\mole}.$
\end{enumerate}
\end{corrige}

\begin{corrige}[attention au piège !]
\begin{enumerate}[label=\alph*)]
  \item En quantité brute, c'est \ce{H2} qui domine (\qty{3}{\mole} $>$ \qty{2}{\mole} de \ce{N2}).
  \item \ce{N2} : $\frac{2}{1}=2$ ; \ce{H2} : $\frac{3}{3}=1$.
  \item Plus petit rapport : \ce{H2} $\Rightarrow$ \textbf{\ce{H2} est limitant}, bien qu'il soit
        \emph{plus abondant} en moles. « Le plus de moles » ne signifie donc pas « en excès » : seul
        le rapport $n/\nu$ décide.
\end{enumerate}
\end{corrige}

\begin{corrige}[précipitation de l'iodure de plomb]
\begin{enumerate}[label=\alph*)]
  \item \ce{Pb(NO3)2} : $\frac{0,1}{1}=0,1$ ; \ce{KI} : $\frac{0,3}{2}=0,15$.
        \textbf{\ce{Pb(NO3)2} est limitant}.
  \item $n(\ce{PbI2})=n(\ce{Pb(NO3)2})\times\dfrac{1}{1}=\qty{0.1}{\mole}.$
  \item \ce{KI} consommé $=0,1\times2=\qty{0.2}{\mole}$, donc restant $=0,3-0,2=\qty{0.1}{\mole}.$
\end{enumerate}
\end{corrige}

\end{document}
